\unchapter{Introduction}

\textbf{Topicality of the problem.}
The rapid advancement of mobile technology has fundamentally transformed how people interact with services in their daily lives.
In the hospitality industry, this transformation is particularly evident in the way customers discover, evaluate, and book restaurants and dining establishments.
The global restaurant technology market reached approximately 4.5~billion~USD in 2023 and is projected to grow at a compound annual growth rate of over 15\% through 2028~\cite{gartner2023}.
In Uzbekistan, traditional teahouses known as ``choyxonas'' embody centuries of cultural heritage and social interaction.
Despite over 70\% smartphone penetration, the restaurant and teahouse booking sector remains largely underserved by modern digital solutions.
Most international platforms such as OpenTable, Resy, and TheFork do not operate in the region, and local establishments continue to rely on phone calls and walk-in visits for reservations.
This gap between growing digital expectations and the absence of a localized booking solution presents both a challenge and a significant opportunity for innovation.

\textbf{Aim of the study.}
The aim of this study is to design, develop, and deploy a comprehensive cross-platform mobile application for automating the reservation process at traditional restaurants and teahouses in Uzbekistan, thereby bridging the gap between established cultural hospitality traditions and modern digital convenience.

\textbf{Object of the study.}
Mobile application development for the hospitality industry in Uzbekistan.

\textbf{Subject of the study.}
Automation of restaurant and teahouse booking processes through a cross-platform mobile application built with Flutter and Firebase.

\textbf{Research problem.}
The absence of a dedicated, culturally adapted digital booking platform for Uzbek restaurants and teahouses leads to inefficient manual reservation processes, limited discovery options for customers, and lack of operational analytics for establishment owners.
Generic international platforms do not address the unique requirements of the Uzbek hospitality market, including private room reservations, extended gathering durations, and multilingual support for Uzbek, Russian, and English languages.

\textbf{Hypothesis.}
Implementing a dedicated cross-platform mobile application specifically designed for the Uzbek hospitality market will reduce reservation processing time by at least 50\% compared to traditional phone-based booking, increase customer satisfaction through improved discovery and booking convenience, and provide establishment owners with data-driven operational insights that enhance business efficiency.

\textbf{Tasks of the study:}
\begin{enumerate}
  \item To review theoretical foundations of mobile application development, analyze existing restaurant booking platforms, and justify the selection of Flutter and Firebase as the technology stack (Chapter~1).
  \item To design the system architecture, database structure, user interface, and real-time synchronization mechanisms for the CHOYXONA.UZ application (Chapter~2).
  \item To implement and describe the functional capabilities of the application across customer, administrator, and super administrator roles (Chapter~3).
  \item To conduct multi-level testing, deploy the application to multiple distribution channels, and outline the future development roadmap (Chapter~4).
  \item To evaluate the results against the research hypothesis and formulate conclusions regarding the effectiveness of the implemented solution.
\end{enumerate}

\textbf{Research methods.}
The methodological basis of this study includes:
\begin{itemize}
  \item \textit{Analysis and synthesis} of scientific literature on mobile development, cross-platform frameworks, and backend services.
  \item \textit{Comparative analysis} of existing booking platforms (OpenTable, TheFork, Resy, Yelp) and cross-platform frameworks (Flutter, React Native, Kotlin Multiplatform).
  \item \textit{Iterative software development} methodology incorporating elements of Agile for practical implementation.
  \item \textit{Prototyping and performance benchmarking} to validate architectural decisions.
  \item \textit{User acceptance testing} with beta testers representing target user personas.
\end{itemize}

\textbf{Thesis structure.}
The thesis consists of four chapters.
Chapter~1 reviews the theoretical foundations including mobile application evolution, booking platform analysis, Flutter framework architecture, and Firebase backend services.
Chapter~2 presents the system architecture and technical design of CHOYXONA.UZ including database structure, UI design, real-time synchronization, and geolocation integration.
Chapter~3 details the functional capabilities across all user roles: customer discovery and booking, administrator management tools, super administrator oversight, and the notification system.
Chapter~4 addresses testing strategies, multi-platform deployment to 15~distribution channels, and the future development roadmap.
The thesis concludes with a synthesis of achievements and validation of the research hypothesis.

The mapping of research tasks to thesis chapters is presented in \cref{tab:task-chapter}.

\begin{table}[htbp]
\caption{Mapping of research tasks to thesis chapters}
\label{tab:task-chapter}
\centering
\begin{tabular}{clc}
\toprule
\textbf{Task} & \textbf{Description} & \textbf{Chapter} \\
\midrule
1 & Literature review and technology justification & Chapter~1 \\
2 & System architecture and database design & Chapter~2 \\
3 & Implementation of functional capabilities & Chapter~3 \\
4 & Testing, deployment, and future roadmap & Chapter~4 \\
5 & Evaluation and conclusions & Conclusions \\
\bottomrule
\end{tabular}
\end{table}
